\clearpage
\pdfbookmark[1]{Nachrichten. 1849 Juli 30.}{anzeige_1849_122_124}
\SourceLeaf{122}
\par\noindent\hfill\rule{0.30\linewidth}{0.4pt}\hfill\mbox{}\par
\begin{center}
\textbf{Nachrichten von der K. Gesellschaft der Wissenschaften zu Göttingen. 1849 Juli 30.}
\end{center}

\par\noindent\hfill\rule{0.30\linewidth}{0.4pt}\hfill\mbox{}\par

 Die in der Sitzung der Königl. Gesellschaft der Wissenschaften am 16. Julius von dem Geh. Hofr. \textsc{Gauss} gehaltene Vorlesung bezog sich auf die
\begin{center}
\textit{Theorie der algebraischen Gleichungen.}
\end{center}
 Von der ausführlichen Denkschrift, die in Kurzem im Druck erscheinen wird, geben wir hier nur folgenden Bericht.

Sie besteht aus zwei Abtheilungen, deren jede einem besondern Gegenstande gewidmet ist.

Die erste Abtheilung beschäftigt sich mit dem Grundlehrsatz der Lehre von den algebraischen Gleichungen, nämlich mit der Zerlegbarkeit jeder rationalen bruchfreien algebraischen Function Einer veränderlichen Grösse in Factoren, worüber der Verf. schon im Jahre 1799 eine Denkschrift veröffentlicht hatte unter dem Titel \textit{Demonstratio nova theorematis, omnem functionem algebraicam rationalem} \textit{integram unius variabilis in factores reales primi vel secundi gradus resolvi posse.} Diese Denkschrift hatte einen doppelten Zweck, zuerst, zu zeigen, dass die sämmtlichen bis dahin versuchten Beweise des ausgesprochenen Lehrsatzes ungenügend und illusorisch waren, und zweitens, einen vollkommen strengen neuen Beweis zu geben. Es würde unnöthig sein, auf den ersten Gegenstand noch einmal zurückzukommen: dagegen aber muss bemerkt werden, dass der Verf. jenem ersten neuen Beweise späterhin noch zwei andere hat folgen lassen, die sich im 3. Bande
\SourceLeaf{123}
der \textit{Commentationes recentiores} unserer Societät befinden, und dass ein vierter von \textsc{Cauchy} zuerst aufgestellt ist. Alle diese vier vollkommen strengen Beweise beruhen auf eben so vielen ungleichen Grundlagen, darin aber kommen sie überein, dass sie sämmtlich \textit{unmittelbar} nur das Vorhandensein Eines Factors der betreffenden Function darthun. Der Strenge der Beweise geschieht hiedurch kein Eintrag. Denn es ist klar, dass nach Ablösung dieses einen Factors von der vorgegebenen Function eine ähnliche Function von niederer Ordnung zurückbleibt, auf welche der Lehrsatz aufs neue angewandt werden kann, und dass aus der fortgesetzten Wiederholung des Verfahrens zuletzt eine vollständige Zerlegung der ursprünglichen Function in Factoren der bezeichneten Art hervorgehen wird. Indessen gewinnt offenbar jede Beweisführung eine höhere Vollendung, wenn gezeigt wird, dass sie geeignet ist, das Vorhandensein sämmtlicher einzelnen Factoren auf einmal unmittelbar anschaulich zu machen. Der erste Beweis ist wirklich in diesem Falle, wie in der gedachten Denkschrift (Art. 23) bereits angedeutet, aber nicht weiter ausgeführt war. Die gegenwärtige Abhandlung gibt nun diese Erweiterung, jedoch nicht in der Gestalt eines Zusatzes zu der Denkschrift von 1799, sondern als Bestandtheil einer Umarbeitung derselben. Es wird nämlich der Beweis hier selbstständig von neuem aufgeführt, aber mit einigen wesentlichen Abänderungen im Gange, wodurch eine bedeutende Vereinfachung gewonnen ist. Verändert ist ausserdem auch die ganze Einkleidung, die bei der ältern Schrift absichtlich so gewählt war, dass alle Einmischung imaginärer Grössen gänzlich vermieden werden konnte. Da jedoch die Beweisführung sowohl ihrem \textit{Ursprunge} als ihrem wahren Wesen nach mit der Conception der complexen Grössen im innigsten Zusammenhang steht, was auch, seitdem diese in der Wissenschaft eingebürgert ist, aufmerksamen Lesern jener Denkschrift nicht entgangen ist, so hat Verf. es für angemessen gehalten, jene erste Einkleidung jetzt fallen zu lassen, und die Argumentation auf ihr eigentliches Feld zurück zu versetzen, deren Grundbegriffe jetzt Jedermann geläufig sind, zumal da dadurch nicht bloss grössere Einfachheit sondern auch noch etwas grössere Allgemeinheit gewonnen wird.

Die zweite Abtheilung ist den algebraischen Gleichungen mit drei Gliedern gewidmet. Diese haben das Eigenthümliche, dass von den zur numerischen Auflösung der Gleichungen bestimmten Methoden einige bei jenen einer Geschmeidigkeit und Eleganz fähig werden, von der ihre Anwendung auf Gleichungen von weniger einfacher Gestalt sehr weit entfernt bleibt. Dies gilt namentlich von der
\SourceLeaf{124}
Auflösung durch unendliche Reihen, und von der indirecten Methode. Es scheint daher die Entwicklung dieser Methoden für die Gleichungen von jener Form eine besondere Ausführung um so mehr zu verdienen, da das Vorkommen solcher Gleichungen in der That ein sehr häufiges ist.

Die Auflösung durch unendliche Reihen hat jedoch der Verf. von seinem gegenwärtigen Zweck gänzlich ausgeschlossen, und nur bemerkt, dass für \textit{jede} Wurzel einer solchen Gleichung, sei sie reell oder imaginär, eine convergente und nach einem leicht erkennbaren Gesetz fortschreitende Reihe gefunden werden kann. So schön aber auch diese Auflösungsart in allgemein theoretischer Beziehung ist, so wird man doch, wo es auf wirkliche praktische Anwendung ankommt, den indirecten Methoden in allen den Fällen den Vorzug geben, wo jene Convergenz nicht eine sehr schnelle ist.

Diese indirecten Methoden nun sind der Gegenstand der zweiten Abtheilung. Es handelt sich hier von \textit{zwei} Methoden, denn das Verfahren, welches zur Bestimmung der imaginären Wurzeln erfordert wird, ist ganz verschieden von dem für die reellen Wurzeln anzuwendenden. Von dem letztern hat der Verf. schon bei andern Gelegenheiten ein paar Proben an besondern Fällen gegeben, und dabei zugleich die allgemeine Anwendbarkeit des Verfahrens angedeutet. Obgleich diese Generalisirung durchaus keine Schwierigkeit hat, so hat der Verf. doch geglaubt, dass man gern die gebrauchfertigen Vorschriften zusammengestellt sehen würde, zumal weil die Anzahl der dabei zu unterscheidenden und einzeln zu behandelnden Fälle nicht unbeträchtlich ist.

Die Auffindung der imaginären Wurzeln auf indirectem Wege ist (insofern nicht schon einigermassen angenäherte Werthe anderswoher bekannt sind) deswegen viel schwieriger, als die der reellen, weil jene aus einem unendlichen Gebiete von zwei Dimensionen herausgesucht werden müssen, diese nur aus einem Unendlichen von Einer Dimension. Diese Schwierigkeit lässt sich bei den Gleichungen von drei Gliedern durch einen einfachen sehr nahe liegenden aber wie es scheint sonst bisher noch nicht benutzten Kunstgriff umgehen. Auch für diese Aufgabe sind die zur Auflösung erforderlichen Vorschriften vollständig und in gebrauchfertiger Gestalt mitgetheilt.
\par\noindent\hfill\rule{0.30\linewidth}{0.4pt}\hfill\mbox{}\par

